\section{What changes when the selected residues can move?}
\label{sec:lower-support}

A lower bound for every labeling and a count of the labelings that do especially well answer different questions. The 64-cycle separates them sharply. The following construction guarantees at least
\[
             z_{\rm window}=\frac{27748320734}{274219308421}
                      =0.101190251313\ldots
\]
under the original score comparison, subject to the arithmetic transport established in this study. A second construction reaches \(11\%\) on 15,318 of 16,384 possible normalized sparse configurations, uniformly over their dense completions. The two rules can be combined by using the second when its certified margin is nonnegative and the first otherwise. This makes support selection materially useful on identifiable inputs. The evidence for the last few difficult configurations is quite different.

The distinction matters for a simple reason. A support can remove an entire difference class by deleting the last pair that realizes a shift. Lowering one endpoint from a large positive rate to a tiny positive rate never accomplishes this in the support-level certificate. Once a shift has disappeared, though, the credit still depends on how much mass remains on the odd side of the node. These two operations face different obstacles. The first is a finite question about the geometry of residues. The second must also respect the score's need for even singleton cells. We isolate their effects by recording three quantities for each tested percentage: the number of covered configurations using one sparse anchor, the number after permitting a second anchor, and how often a passing maximizing support uses a sparse mask absent from the historical menu.

\subsection{Retaining the source's whole window budget}

Write \(h=2+8\varepsilon\), \(0<\varepsilon<1/52\), and \(r=r_E(\varepsilon)\). Recall the means \(\mu\) of all sites and \(\nu\) of even sites in one candidate cell. A literal marked tile has limiting independent probability
\[
        \theta=\exp(-(6r+1)\mu)\,\nu^{3r}.
\]
The convenient sufficient inequality \(\theta>10^{-(6r+1)}\) can spend more of the failure budget than necessary. There are at least \((h+3)10^{6r+1}\log x\) disjoint trials after dropping the possible trial containing a conditioned retained point. There are \(O(x^h)\) starts. The weighted score proof asks for a per-orbit conditional failure probability \(o(x^{-2})\). Consequently it is enough that
\begin{equation}
       10^{6r+1}\theta > \frac{h+2}{h+3}.
       \label{eq:lower-window}
\end{equation}
Indeed, strictness permits a \(\theta'<\theta\) on the same side of the inequality. Each late trial has success probability at least \(\theta'\) in the comparison law. Multiplying the exponential failure bound by the number of starts gives the exponent \(h-(h+3)10^{6r+1}\theta'<-2\). The relative fixed-count comparison is applied at each start before the union. Both \(r\) and the strict margin are fixed before \(x\) increases. The derivation uses the literal window multiplicity in the source and applies to the chosen sampling argument. A better estimate of score compatibility for another law could give a different feasible set.
\href{https://github.com/nousentllc/riemann-hypothesis/blob/40e56bab026552e089d7afcfeff08d2defb9c005/src/Formalism/SharpRestriction/GrowthBounds/UnramifiedQuadraticSetting.lean#L239-L251}{Window multiplicity in the original source}.
\href{https://github.com/nousentllc/riemann-hypothesis/blob/40e56bab026552e089d7afcfeff08d2defb9c005/src/Development/AutomorphicArithmetic/ThickFrames/TilesLe.lean#L43-L67}{Trials after conditioning on one site}.

For a numerical certificate uniform over \(\varepsilon\), factor the left side of \eqref{eq:lower-window} exactly:
\[
       10^{6r+1}\theta
       =\frac{(10e^{-\mu}\sqrt{\nu})^{6r+1}}{\sqrt{\nu}} .
\]
Since \((h+2)/(h+3)<54/67\), the two rationally verifiable tests
\[
    \log 10-\mu+\tfrac12\log\nu>0,
    \qquad (54/67)^2\nu<1
\]
suffice on a tile. Their separation is deliberate. The exponent can grow without bound as \(\varepsilon\) decreases. A base strictly below one cannot be offset for all \(\varepsilon\) by a fixed prefactor. Conversely a base slightly above one, accompanied by a suitable prefactor, already controls every finite exponent. This explains both the gain and its small size: the asymptotic leaf constraint itself moves very little.

In the exceptional frozen support give 26 ordinary high-rate leaves \(c=2.392096615\), the even singleton of the sparse half \(b=0.706267\), and the minimum of the six-leaf dense-quarter pattern \(a=64-26c-b=1.09922101\). There are 23 distinct pairs of linear forms for the two means across all three permissible reservoir deletions and all depths. Twenty-two satisfy the ordinary strict product test, checked with rational logarithm enclosures. The remaining pair is precisely \((c,c/2)\). Both inequalities above hold for this pair. Every rate is positive and less than \(12/5\), so the separate upper-count estimate used in the score transfer continues to apply. The finite geometry from the fixed-support analysis gives
\[
                 \frac PQ\ge\frac{116c}{2112+254c+32b}
                    =z_{\rm window}.
\]
The ordinary charts give at least \(11661/115010\), which exceeds the displayed floor. The rational parameters have deliberately retained strict margins. This certificate provides a counterexample to treating the optimum of the standardized tile test as the optimum for the literal source score. Its improvement over the strongest frozen standardized certificate is about \(0.00003555\) percentage points. The preceding section approaches the same literal corner more closely with sharper rational rates. The wider margins here provide a shorter independent check before the support experiment.

\subsection{An exhaustive sparse-coordinate experiment with explicit quantifiers}

Denote by \(E,O\) the preferred and opposite halves of the root. Decompose their children in canonical order as \((A,B)\) and \((C,D)\), respectively; each has width 16. The letters in this section always refer to these four quarters. In \(A,B,D\) permit every nonempty canonical mask satisfying the local even-mass tests described below. In \(C\), retain its own minimum and optionally the minimum in its opposite child. Each positive scaled rate belongs to
\[
            \{1/8,2/8,\ldots,19/8\}.
\]
The total is 512 eighths. This finer control of the locations and coarser choice of the rates are independent design choices. In particular, this family can fail just below the continuous fixed-support certificate. Falling back on that certificate keeps the all-label result.

Give \(O\) coordinates in \(\mathbb Z/32\) with the \(C\) minimum at zero and multiply by the inverse of the coordinate of the \(D\) minimum to put it at one. Its oddness makes the inverse available. The selected \(C\) coordinates are \(2S\), where either \(S=\{0\}\) or \(S=\{0,s\}\), with \(s\) odd modulo 16. The \(D\) coordinates are \(1+2T\). The recursive canonical order on a 16-cycle comes from 15 binary choices. Of the resulting \(2^{15}\) orders, exactly 2,048 have their first entry at zero, the condition on the coordinate list of \(D\). Every odd \(s\) can occur separately for each such order since the two subtrees of \(O\) have separate choices. Hence the experiment tests \(2{,}048\cdot8=16{,}384\) sparse configurations.

Within either dense quarter, shift its first canonical point to zero and multiply the odd difference to its opposite-child minimum by its inverse. The conditions on canonical entries 0 and 8 reduce the number of orders to \(2^{15}/(16\cdot8)=256\). The affine change preserves addition of steps, the number of missing steps at each depth, and canonical odd mass. Different dense quarters may choose different normalizations because the conservative search takes zero recovery credit at the width-64 root and at the width-32 node \(E\). It still charges their entire odd exposure. This is what makes a bound for the two independently normalized quarters applicable to every completion in the physical cycle. To credit root recovery one would also have to align their actual relative coordinates.

Let \(T\) now denote the occupied coordinate set within \(D\). The digits \(j\) corresponding to the allowed odd increments \(1+2j\) in \(O\) are exactly
\[
                   (T-S)\ \cup\ \{-1-k:k\in T-S\}.
\]
Taking differences from a point of \(2S\) to a point of \(1+2T\) gives the first set. Taking differences in the opposite direction gives the second. Thus the number \(d_O\) of missing increments is 16 minus the size of that union. The second anchor can enlarge the union, which reduces the return on every tick in \(D\). Locally inside \(C\) the two anchors have opposite parity. They realize exactly two distinct odd differences, \(s\) and \(-s\); the other six odd differences in that 16-cycle are missing. The extra gain for mass on the second anchor is therefore 12 times its mass. If the two new translated difference sets contribute no previously uncovered digit to the union, the additional anchor earns this local gain with no reduction in \(d_O\). If they cover several new digits, a singleton can be preferable. The search takes both possibilities seriously.

\subsection{A threshold computation that preserves the signs}

At target \(z=p/q\) write the integer margin as \(qP_8-pQ_8\), where the subscript denotes rates in eighths. Clearing denominators lets all geometry and all threshold decisions use integers. Inside a quarter the coefficient of a single tick at canonical rank \(i\) is
\[
 2q\!\sum_{\substack{v\text{ internal to the quarter}\\i\text{ in its odd child}}}\!d_v
                     -p(1+2i).
\]
Indeed, the weights of the binary digits of rank \(i\), doubled, are exactly the widths of the internal nodes where its branch was odd; the extra one accounts for the leaf tail. Every \(B\) tick pays another \(32p\), and every \(D\) tick adds \(2q d_O-96p\). The first \(C\) anchor contributes \(-65p\) per tick and the second contributes \(12q-81p\). In particular dropping parent recovery never removes its exposure from the denominator.

To obtain a reusable finite constraint on marks, take an occupied local node of width \(w\), \(n_v\) active leaves and total mass at most \(19n_v\) ticks. Use the upper mean \(19n_v/(8w)\) in the factor \(e^{-\mu}\). Require the even-child tick count \(u\) to be at least the smallest integer for which
\[
         e^{-19n_v/(8w)}
         \min\left\{\left(\frac{u}{8w}\right)^{1/2},
                    \left(\frac{u}{8w}\right)^{6/13}\right\}>\frac1{10}.
\]
The two exponents bound the source interval since \(r\ge2\). For a given width and total, the checker encloses the exponential by a Taylor partial sum and a geometric upper bound on the positive remainder. It verifies that the first passing integer clears the upper endpoint and its predecessor fails even at the lower endpoint. The code also checks every homogeneous leaf rate; otherwise all coarse nodes could pass while the increasingly fine levels fail. In the two upper levels and in the sparsely occupied \(C\) quarter it checks the mark again at the actual masses. Exactly 1,806 of the 65,535 nonempty quarter masks survive the conservative local test.

This organization makes changes in mass cheap to evaluate. Give an active leaf one tick to retain its support. Visit the mandatory even-subtree floors from the smallest nodes upward. To discharge a deficit add available ticks in that subtree in descending order of margin coefficient. A later addition can only help earlier lower bounds, since the assumed total in the exponent was already at the support maximum. After all floors pass, fill remaining capacity in descending coefficient order and record one feasible objective at every total reached. The certificate needs only the feasibility and the achieved objective of these allocations. Stronger constraints from the worst total may discard arrays whose actual marks would pass. In particular a negative outcome of this routine cannot upper-bound all profiles on its chosen support.

For a dense order \(L\), let \(V_L(m)\) be the best margin offered at total \(m\) in its quarter. Take the pointwise lower envelope \(V(m)=\min_L V_L(m)\). At given combined dense total \(e\), every dense completion can guarantee
\[
  W(e)=\max_{a+b=e}\{V(a)+V(b)-32pb\},
\]
with the mark required at \(E\). A maximizing split \((a,b)\) is universal at that total, while the two winning internal supports may depend on their own orders. Combining \(W\) with each sparse configuration and both choices of \(S\) yields a lower bound before reading the dense part of the labeling. The finite operations are in this order because choosing a different favorable dense order for each mass would answer the wrong quantifier. The minimum inside the envelope makes the certificate conservative and suggests a precise check when a sparse configuration misses its target.

\subsection{Where the additional support pays for itself}

The following counts refer to the normalized \((D,s)\) coordinate pairs just defined. A passing count applies to every dense completion of that pair. The last column counts passing cases in which the chosen maximum uses a \(D\) mask outside the union of all eleven masks in the preceding decision rule. Ties are resolved by the enumeration order. The count identifies actual use of a new mask; it does not establish that no historical mask could also have cleared the same threshold after the other quarters changed.
\begin{center}
\begin{tabular}{@{}rrrrr@{}}
\toprule
Target & Single anchor & Either anchor rule & Additional passes & New \(D\) mask\\
\midrule
\(10.12\%\) & 16,336 & 16,369 & 33 & 14,035\\
\(10.30\%\) & 16,328 & 16,364 & 36 & 14,018\\
\(10.50\%\) & 16,288 & 16,328 & 40 & 14,001\\
\(11.00\%\) & 15,152 & 15,318 & 166 & 13,585\\
\bottomrule
\end{tabular}
\end{center}
At \(11\%\), the two-anchor option rescues 166 cases that the single-anchor component-envelope search could not certify at the same threshold. For one such layout the literal physical witness uses the following rate ticks. In \(A\) take rate 19 on all 16 leaves. In \(B\) give rank 0 rate 16, and ranks \(8,10,11,12,13\) rate 19. In \(C\) take rates 2 and 19 on ranks 0 and 8. In \(D\) give ranks \(0,4,8,12\) rate 19. The four quarter totals are \(304,111,21,76\), summing to 512. The positive \(B\) ranks in the saved physical layout correspond respectively to \(2,6,22,54,14,46\) modulo 64. The positive \(C\) and \(D\) coordinates are \(1,5\) and \(3,11,7,15\).

As a hand check, the two normalized sets are \(S=\{0,1\}\) and \(T=\{0,1,2,3\}\). Their difference digits run from \(-1\) to 3. After adjoining the reflected digits they run from \(-4\) to 3 modulo 16. The other eight digits give the missing odd increments in \(O\). Thus its recovery is \(2\cdot8\cdot76=1216\) ticks. The next level contributes \(2\cdot2\cdot95\) in \(B\), \(2\cdot6\cdot19\) in \(C\), and \(2\cdot4\cdot38\) in \(D\), or 912 in total. Two width-eight nodes contribute 76 each. The global root and dense parent have no blocked steps. Detailed allowance charges can be checked without summing individual leaf coefficients:
\begin{center}
\begin{tabular}{@{}lrrrrrrr@{}}
\toprule
Contribution & 64 & 32 & 16 & 8 & 4 & 2 & Tail\\
\midrule
Recovered ticks & 0 & 1216 & 912 & 152 & 0 & 0 & 0\\
Allowance ticks & 6208 & 5984 & 4864 & 1824 & 760 & 380 & 512\\
\bottomrule
\end{tabular}
\end{center}
Their ratio is \(2280/20532=190/1711=11.104617\ldots\%\). An independent replay reads the physical coordinates in the witness, subtracts each possible odd step at every occupied node, checks the actual quantitative marks, and agrees with the threshold search on its integer margin \(100(2280)-11(20532)=2148\). This example also illustrates why root recovery alone is a poor measure of a profile's quality. Here three smaller widths carry the full numerator while every width contributes to exposure.

\subsection{Taking apart the fifteen failures}

At \(10.12\%\), fifteen of the 16,384 sparse configurations fail the pointwise lower-envelope test. To distinguish a difficult geometry from loss through the order of the extrema, retain the entire curve \(V_L\) of each dense order. For each missed sparse configuration, first compute \(J(e)\), the best sparse contribution at residual dense total \(e\) over the same admissible one- and two-anchor choices. Combine it with every ordered pair of dense orders, now testing the sign of
\[
 \max_{a,b}\{V_{L_A}(a)+V_{L_B}(b)-32pb+J(a+b)\}.
\]
The actual mark at \(E\) is required for each tested split. This allows the split to depend on the dense orders. Pointwise-identical payoff curves can be processed together with their exact multiplicities. The 256 normalized orders give 121 distinct curves at this threshold, and the sum of products of multiplicities is exactly \(256^2=65{,}536\).

The complete Cartesian computation recovers 65,280 of the 65,536 dense completions for \emph{each} of the fifteen former misses. Thus this two-anchor family covers at least \(2^{30}-3{,}840\) of the \(16{,}384\cdot256^2=2^{30}\) normalized joint chart tuples at \(10.12\%\). This count carries more information than a further decimal place on a uniform guarantee. Nearly all the apparent failures resulted from making the dense total insensitive to the labels. The 256 remaining ordered pairs in each row all have the same \(B\) order and permit every one of the 256 orders on \(A\). The sole normalized list on \(B\), number 11,776 in the enumeration, is \((0,8,12,4,14,6,10,2,1,9,5,13,3,11,7,15)\). This locates the geometric obstruction in a specific dense quarter, even after allowing each dense order its own split.

The fallback from the fixed-support experiment localizes the survivors further. Reflect \(D\)'s first eight canonical physical coordinates about the \(C\) minimum. In normalized coordinates this sends \(j\) to \(-1-j\bmod16\). Consult the historical rule unless this reflection pairs the first eight canonical ranks with the last eight in reverse order. In the reversed case write the first eight coordinates as \(2\operatorname{rev}_3(g(i))\). The historical exceptional decision requires its free switch bits at indices \((2,4,5,6)\) to be \((0,1,1,0)\). The other bits are fixed because \(g(0)=0\). Reversing the three bits of each coordinate reproduces
\[
    g=(0,1,3,2,5,4,6,7).
\]
Hence the sole exceptional sixteen-coordinate list is
\[
 (0,8,12,4,10,2,6,14,1,9,13,5,11,3,7,15).
\]
An independent pass through all 2,048 normalized lists confirms exactly one such list. It is order 10,320 in the saved enumeration of all width-16 orders. This classification uses only the \(D\) order relative to the first \(C\) anchor. Neither the position of the permitted second anchor nor the two independently normalized dense quarters affects the historical case split. Of the fifteen residual sparse configurations only three have this \(D\) list, with \(s=3,5,7\). In the other twelve the ordinary frozen guarantee already exceeds \(10.12\%\).

It follows that the combined certificate reaches \(10.12\%\) on at least \(2^{30}-3\cdot256=1{,}073{,}741{,}056\) normalized tuples. Every realization of a passing tuple enjoys the guarantee. The count expresses coverage of a combinatorial classification. Arithmetic label frequencies and independence are unnecessary for the assertion and are unknown here. On the remaining 768 tuples the fallback still gives the all-label bound. A matching upper bound for these residual tuples would have to allow alternative \(C\) supports, broader rates and the physical alignments discarded by the two independent dense normalizations. The exact count of residual tuples supplies a finite target for that work.


We ran one more targeted test on these residuals. Keep \(B\) at order 11,776 and \(D\) at 10,320. For each of the three steps read the complete \(C\) order, of which there are 256 possibilities after fixing its two marked coordinates. Permit every structurally compatible \(C\) mask with up to eight leaves and every such \(D\) mask with up to nine. Give all active leaves 19 ticks except each quarter's minimum, independently allowed every positive tick count through 19. For each candidate subtract the actual physical differences in the 32-cycle \(O\), recheck local marks at the actual rates, and combine the result with the known \(B\) order and the universal dense-\(A\) curve. All 768 detailed sparse configurations returned the same best integer margin, \(-136{,}968\) for target \(1012/10000\). Here the 768 configurations index three steps and 256 \emph{internal C orders}; the preceding residual count of 768 indexed those steps and 256 \emph{A orders}. The two counts describe different tests. This failed extension identifies a reason to change the next search: allowing more positive leaves while tying all their non-minimum rates to the cap did not recover a single additional configuration. For a particular surviving chart, interior rates or additional root credit could still help. The adverse-word calculation later shows why such per-chart improvements cannot certify 10.12 percent uniformly within the marked family.
